% ===================================================================
%  LENTELĖ Nr. 6 — Namas iš vidaus, baldai, maitinimas, apšvietimas.
%  Traduction 2026 d'apres texto/io, controlee sur texto/fr, texto/ru
%  et texto/cs. Consigne : voir texto/lt/10-lentele-01.tex.
%
%  LE PLUS LONG TABLEAU DU LIVRET, cinq pieces et plus de deux cents
%  objets, et le second banc d'essai de la colonne apres le corps du
%  tableau 2.
%
%  IL DONNE UN RESULTAT QUE LES TREIZE COLONNES PRECEDENTES NE
%  POUVAIENT PAS DONNER : ici, LE NOM DIT L'ACTION. Le lituanien
%  nomme l'ustensile par un suffixe pose sur le verbe, et l'objet
%  porte donc son emploi dans son nom :
%
%    « žadintuvas » ce qui reveille (le reveil),
%    « laistytuvas » ce qui arrose (l'arrosoir),
%    « šildytuvas » ce qui chauffe (le chauffe-bain),
%    « degtukas » le petit brulant (l'allumette),
%    « kavamalė » ce qui moud le cafe,
%    « skambutis » le petit sonnant (la sonnette),
%    « keptuvė » ce qui frit (la poele),
%    « virtuvė » l'endroit ou l'on cuit (la cuisine),
%    « skalbykla » l'endroit ou l'on lave,
%    « prausykla » l'endroit ou l'on se debarbouille.
%
%  Et les pieces d'habitation sont des participes : « valgomasis »
%  celui ou l'on mange, « miegamasis » celui ou l'on dort — deja vus
%  au tableau 5.
%
%  CE N'EST PAS UN EMPRUNT REFUSE, C'EST UN EMPRUNT INUTILE. Le
%  tcheque, treizieme colonne, avait du DECIDER contre l'allemand ;
%  le lituanien n'a pas eu a decider, parce que la langue fabrique
%  elle-meme, avec deux ou trois suffixes vivants, tout le
%  vocabulaire d'une maison. Le tableau 5 avait montre la limite de
%  ce pouvoir — la maconnerie, arrivee du dehors, a garde ses mots ;
%  le tableau 6 en montre la portee.
%
%  ET CELA REPOND A LA NOTE DU FEUILLET 17, celle des jeux, qui
%  autorisait l'ido a nommer soit par l'action, soit par le nom
%  national. Ce que l'ido s'accordait comme un expedient, le
%  lituanien le fait par defaut.
%
%  L'ECART DECLARE DU BLOC c1-06-1, LE PREMIER DES TROIS. Le
%  fac-simile ido y ouvre sur « (6) », qui est le savon de l'alinea
%  2 ; c'est la femme de chambre « (16) », comme le porte le francais
%  et comme la planche le grave. La colonne le rend tel quel, avec
%  les treize qui la precedent, et renvoji.py le passe sans rien
%  dire.
% ===================================================================

%%K t06-apar-1 apar
\VUsaut{6.0mm}
\VUtitre{15.0pt}{60}{LENTELĖ Nr. 6}
\VUsaut{1.4mm}
\VUtitre{11.4pt}{0}{Namas iš vidaus, baldai, maitinimas,}
\VUsaut{0.4mm}
\VUtitre{11.4pt}{0}{apšvietimas.}
\VUsaut{2.2mm}

%%K t06-apar-2 apar
Namas baigtas. Jono dėdė įsikūrė savo naujoje buveinėje, kurios išplanavimą žinome.
Dabar pažiūrėkime, kaip jis apstatė savo namą. Pirmiausia apžiūrėsime:

%%K t06-tit-1 sub
\VUpk{10.2pt}{40}{Miegamoji.}
\VUsaut{3.0mm}

%%K t06-c1-01-1 p
1. --- Atvykstame nelaiku, nes kambarinė užsiėmusi valydama \VUgras{kambarį}.

%%K t06-c1-02-1 p
2. --- Ant \VUgras{prausimosi stalelio} \textsuperscript{(1)} šen ir ten guli įvairūs
daiktai, kuriais prausiamasi, šukuojamasi, trumpai tariant, tvarkomasi. Tai
\VUgras{dubenėlis} \textsuperscript{(2)} ir \VUgras{ąsotis} \textsuperscript{(3)},
\VUgras{galvos šepetys} \textsuperscript{(4)}, \VUgras{šukos} \textsuperscript{(5)},
\VUgras{muilas} \textsuperscript{(6)} ir \VUgras{kempinė} \textsuperscript{(7)},
galiausiai \VUgras{buteliukas} \textsuperscript{(8)} skystiems dantų vaistams.

%%K t06-c1-03-1 p
3. --- Ant grindų, prieš prausimosi stalelį, štai \VUgras{kibiras}
\textsuperscript{(9)} nešvariam vandeniui surinkti.

%%K t06-c1-04-1 p
4. --- \VUgras{Prieškambaryje} \textsuperscript{(10)} matome kelias \VUgras{pakabas}
\textsuperscript{(13)} ir du \VUgras{skėčius} \textsuperscript{(11)}
\VUgras{skėtinėje} \textsuperscript{(12)}.

%%K t06-c1-05-1 p
5. --- Kitoje durų pusėje yra \VUgras{spinta su veidrodžiu} \textsuperscript{(14)},
kurioje laikomi \VUgras{skalbiniai} \textsuperscript{(15)}.

%%K t06-c1-06-1 p
6. --- Pasinaudokime tuo, kad \VUgras{kambarinė} \textsuperscript{(6)} kloja
\VUgras{lovą} \textsuperscript{(17)}, ir apžiūrėkime įvairias jos dalis. \VUgras{Lovos
rėmas} \textsuperscript{(20)} laiko gerą \VUgras{spyruoklinį čiužinį}
\textsuperscript{(21)}, ant kurio Justina tuoj padės vilnonį \VUgras{čiužinį}
\textsuperscript{(22)}. Paskui ji paims nuo kėdžių dvi \VUgras{paklodes}
\textsuperscript{(23)} ir \VUgras{antklodę} \textsuperscript{(24)}. Tada ji padės
lovos galvūgalyje \VUgras{skersinę pagalvę} \textsuperscript{(25)} ir
\VUgras{pagalvėlę} \textsuperscript{(26)}, kurios guli su \VUgras{pūkine antklode}
\textsuperscript{(27)} ant \VUgras{lovos kilimo} \textsuperscript{(32)}. Plunksnų
šluostikliu ji nudulkina \VUgras{užuolaidas} \textsuperscript{(19)} ir \VUgras{lovos
baldakimą} \textsuperscript{(18)}, ir štai dalis darbo atlikta.

%%K t06-c1-07-1 p
7. --- Tada jai reikės kiek sutvarkyti \VUgras{lovos staliuką}
\textsuperscript{(28)}, prisukti \VUgras{žadintuvą} \textsuperscript{(29)}, paruošti
\VUgras{naktinę lempelę} \textsuperscript{(30)}, išimti \VUgras{žvakę}
\textsuperscript{(31)}, kad išvalytų žvakidę.

%%K t06-tit-2 sub
\VUsaut{5.0mm}
\VUpk{10.2pt}{40}{Darbo krepšys ir židinio reikmenys.}
\VUsaut{3.0mm}

%%K t06-c1-08-1 p
8. --- \VUgras{Darbo krepšys} \textsuperscript{(34)} prie \VUgras{taburetės}
\textsuperscript{(33)} taip pat labai reikalingas sutvarkyti. Jai reikės sulankstyti
\VUgras{siuvinius} \textsuperscript{(35)}, sudėti į \VUgras{darbo staliuką}
\textsuperscript{(38)} \VUgras{antpirštį}, \VUgras{siūlus} \textsuperscript{(36)},
\VUgras{adatas} \textsuperscript{(37)} ir \VUgras{žirkles} \textsuperscript{(39)} ir
\VUgras{raktu} \textsuperscript{(40)} užrakinti staliuko dangtį.

%%K t06-c1-09-1 p
9. --- Ant \VUgras{vienkojo staliuko} \textsuperscript{(41)} viduryje Justina
pakeis \VUgras{grafino} \textsuperscript{(42)} vandenį. Ji įpils \VUgras{cukraus} į
\VUgras{cukrinę} \textsuperscript{(43)}, nuvalys \VUgras{cukraus žnyples}
\textsuperscript{(44)} ir padės \VUgras{drabužių šepetį} \textsuperscript{(46)}.

%%K t06-c1-10-1 p
10. --- Dešinioji kambario pusė pareikalaus iš jos mažiau darbo, nes \VUgras{ilgas
krėslas} \textsuperscript{(47)} ir jo \VUgras{pagalvė} \textsuperscript{(48)} yra savo
vietoje, o kadangi oras nešaltas, niekas nepajudinta tarp \VUgras{židinio}
\textsuperscript{(49)} reikmenų (padargų). \VUgras{Semtuvėlis}
\textsuperscript{(50)}, \VUgras{žnyplės} \textsuperscript{(51)}, \VUgras{pagaikščiai}
\textsuperscript{(55)}, \VUgras{pelenų užtvarėlė} \textsuperscript{(56)} švarūs;
\VUgras{ugnies ekranas} \textsuperscript{(54)} nebuvo pajudintas, o \VUgras{malkinė}
\textsuperscript{(52)} gerai aprūpinta \VUgras{malkomis} \textsuperscript{(53)}.

%%K t06-noto-1 noto
\VUnotes{143.2mm}{%
(*) Babo = mažas vaikas; angl. \textit{baby}, pranc. \textit{bébé}, it.
\textit{bambino}.%
}

%%K t06-noto-2 noto
\VUnotes{143.2mm}{%
(*) Lambrequino = pranc. \textit{lambrequin}, isp. \textit{lambrequines}, port.
\textit{lambrequins}, net vok. ir angl. \textit{lambrequins}, taigi vok. angl.
pranc. port. isp.%
}

%%K t06-c1-11-1 p
11. --- Vis dėlto jai reikės nudulkinti \VUgras{švytuoklinį laikrodį}
\textsuperscript{(57)}, \VUgras{žvakides} \textsuperscript{(58)} ir \VUgras{smulkmenų
dėžutes} \textsuperscript{(59)}, galiausiai nušluostyti \VUgras{veidrodį}
\textsuperscript{(60)}.

%%K t06-c1-12-1 p
12. --- O kūdikio (babo (*)) \VUgras{lopšys} \textsuperscript{(61)} jau paruoštas.

%%K t06-tit-3 sub
\VUsaut{5.0mm}
\VUpk{10.2pt}{40}{Svetainė.}
\VUsaut{3.0mm}

%%K t06-c2-01-1 p
1. --- Dabar štai svetainė. Tai labai gražus \VUgras{kambarys.} Židinys, kuris
stovi dešiniajame kampe, papuoštas dailia \VUgras{statulėle} \textsuperscript{(1)},
dviem gražiomis \VUgras{vazomis} \textsuperscript{(3)} ir apdrapiruotas aksominiu
\VUgras{lambrekinu} \textsuperscript{(2)} (*).

%%K t06-c2-02-1 p
2. --- Kad židinio žiežirbos neuždegtų kilimo, prieš židinį pastatytas varinis
\VUgras{žiežirbų sietelis} \textsuperscript{(4)}.

%%K t06-c2-03-1 p
3. --- Netoliese atvertas elegantiškas moteriškas \VUgras{rašomasis staliukas}
\textsuperscript{(5)}. Prie jo \VUgras{namų šeimininkė} \textsuperscript{(23)} rašo
savo laiškus.

%%K t06-c2-04-1 p
4. --- Langas papuoštas \VUgras{lango užuolaidomis} \textsuperscript{(6)} su
\VUgras{raiščiais} \textsuperscript{(8)}, užkabintais ant \VUgras{kabliukų}
\textsuperscript{(7)}, ir medžiaginėmis portjeromis, \VUgras{viršuje
apdrapiruotomis} \textsuperscript{(9)}.

%%K t06-c2-05-1 p
5. --- Lango nišoje \VUgras{vazone} \textsuperscript{(11)} auga žalias
\VUgras{augalas} \textsuperscript{(10)}, puiki palmė, kurios lapai driekiasi
beveik iki \VUgras{žibalinės lempos} \textsuperscript{(12)}.

%%K t06-c2-06-1 p
6. --- Lempos \VUgras{gaubtą} \textsuperscript{(13)} sutvarkė Marija, žavinga
mergina \textsuperscript{(14)}, kuri sėdi prie pianino \textsuperscript{(15)}. Labai
tiesiai sėdėdama ant \VUgras{taburetės} \textsuperscript{(16)}, ji mokosi muzikos
kūrinį. Ji labai sumani ir rūpestinga, kaip rodo jos \VUgras{natų spintelė}
\textsuperscript{(17)}, kuri puikiai sutvarkyta.

%%K t06-tit-4 sub
\VUsaut{5.0mm}
\VUpk{10.2pt}{40}{Padauža.}
\VUsaut{3.0mm}

%%K t06-c2-07-1 p
7. --- Jos brolis Paulius ieško \VUgras{knygos} \textsuperscript{(19)}
\VUgras{knygų spintoje} \textsuperscript{(18)}. Jis \VUgras{jaunuolis}
\textsuperscript{(20)}, judrus ir nepakenčiamas. Kabindamasis ant spintos, kad
pasiektų knygą, kuri per aukštai, jis rizikuoja nuversti \VUgras{biustą}
\textsuperscript{(21)}, esantį viršuje.

%%K t06-c2-08-1 p
8. --- Tai kvailystei jis pasinaudoja tuo, kad motina užsiėmusi kalbėdama su
drauge, \VUgras{ponia} \textsuperscript{(22)}, kuri atvyko praleisti keleto dienų jos
namuose. Motina sėdi ant \VUgras{sofos} \textsuperscript{(24)} prie \VUgras{krėslo}
\textsuperscript{(25)}, o jos draugė ant \VUgras{kėdės} \textsuperscript{(27)},
kurios matyti tik \VUgras{atlošas} \textsuperscript{(26)}.

%%K t06-c2-09-1 p
9. --- Didelis \VUgras{rėmas} \textsuperscript{(30)}, kabantis ant sienos, laiko
giminaitės \VUgras{portretą} \textsuperscript{(29)}. \VUgras{Albume}
\textsuperscript{(32)}, padėtame ant stalo, surinktos kitų šeimos narių
\VUgras{nuotraukos} \textsuperscript{(33)}. Vaikai ką tik jį vartė, nes
\VUgras{kėdės} \textsuperscript{(31)} dar sustumtos aplink stalą.

%%K t06-c2-10-1 p
10. --- Porcelianinė \VUgras{kiniška vaza} \textsuperscript{(34)}, kuri stovi prie
\VUgras{popieriaus peilio} \textsuperscript{(35)}, buvo padovanota Marijai per jos
šventę, o \VUgras{širmą} \textsuperscript{(36)}, kuri puošia kambario kampą, iš
Kinijos parvežė jos dėdė.

%%K t06-c2-11-1 p
11. --- Pralinksminta puikių \VUgras{paveikslų} \textsuperscript{(37)}, kabančių ant
\VUgras{sienos} \textsuperscript{(42)}, apšviesta \VUgras{lubinės liustros}
\textsuperscript{(38)}, kurios \VUgras{elektros lemputės} \textsuperscript{(39)}
vakare linksmai spindi, ta svetainė atrodo labai maloni. Minkštas \VUgras{kilimas}
\textsuperscript{(40)}, patiestas ant parketo, ir toks patogus \VUgras{elektrinis
skambutis} \textsuperscript{(41)} užbaigia jos jaukumą.

%%K t06-tit-5 sub
\VUsaut{3.6mm}
\VUpk{10.2pt}{40}{Valgomasis.}
\VUsaut{3.0mm}

%%K t06-c3-01-1 p
1. --- Suskambo vakarienės varpas. Visa šeima, išskyrus Mariją, susirinkusi aplink
stalą. Valgomasis maloniai šildomas puikios fajansinės \VUgras{krosnies}
\textsuperscript{(1)} su plonu \VUgras{vamzdžiu} \textsuperscript{(2)}.

%%K t06-c3-02-1 p
2. --- Tame kambaryje nedaug baldų. Palei sieną kabo, virš \VUgras{suoliuko}
\textsuperscript{(4)}, puošnus \VUgras{fontanėlis} \textsuperscript{(3)}. Visai
netoli yra \VUgras{indauja} \textsuperscript{(5)}, ant kurios dedami šalti valgiai,
deserto lėkštės ir įvairūs stalo reikmenys, kurių tuoj pat nereikia.

%%K t06-c3-03-1 p
3. --- Taip ant viršutinės lentynos matome \VUgras{keptą viščiuką}
\textsuperscript{(7)}, \VUgras{žuvį} \textsuperscript{(8)} ir \VUgras{dubenį}
\textsuperscript{(10)} su \VUgras{vaisiais} \textsuperscript{(9)}. Žemiau štai
\VUgras{sūris} \textsuperscript{(11)}, \VUgras{duona} \textsuperscript{(12)},
\VUgras{prieskonių rinkinys} \textsuperscript{(13)}, kuriame yra visa, ko reikia
salotoms pagardinti: aliejus, actas, \VUgras{druska} \textsuperscript{(14)} ir
\VUgras{pipirai} \textsuperscript{(15)}. Galiausiai ant apatinės lentynos yra
\VUgras{pyragas} \textsuperscript{(17)} prie \VUgras{garstyčinės}
\textsuperscript{(16)}.

%%K t06-c3-04-1 p
4. --- Valgomojo laikrodis, vadinamas \VUgras{sieniniu laikrodžiu}
\textsuperscript{(20)}, labai savitos formos. Jis įtaisytas mediniame rėme, kuriame
dar yra \VUgras{barometras} \textsuperscript{(19)} ir \VUgras{termometras}
\textsuperscript{(18)}.

%%K t06-tit-6 sub
\VUsaut{4.6mm}
\VUpk{10.2pt}{40}{Vakarienės patarnavimas.}
\VUsaut{3.0mm}

%%K t06-c3-05-1 p
5. --- Prie visų kambario sienų pritvirtintos vaškuoto ąžuolo \VUgras{medinės
plokštės} \textsuperscript{(21)}. Medžiaginė \VUgras{durų portjera}
\textsuperscript{(22)} pakankamai uždaro duris, pro kurias einama į verandą. Ten
šeima eis gerti kavos po vakarienės. Jau \VUgras{puodeliai} \textsuperscript{(24)} ir
\VUgras{lėkštutės} \textsuperscript{(25)} paruošti ant \VUgras{indaujos}
\textsuperscript{(23)}.

%%K t06-c3-06-1 p
6. --- Virš stalo prie lubų pritvirtintas \VUgras{pakabinamas šviestuvas}
\textsuperscript{(26)}, kurio labai ryški lempa vakare apšviečia visą valgomąjį.

%%K t06-c3-07-1 p
7. --- Dabar apžiūrėkime patį \VUgras{valgomąjį stalą} \textsuperscript{(27)}. Kai
šeima įėjo, indai buvo paruošti. Ant \VUgras{staltiesės} \textsuperscript{(28)} prie
kiekvieno stalininko vietos buvo padėta \VUgras{lėkštė} \textsuperscript{(33)},
\VUgras{servetėlė} \textsuperscript{(29)}, \VUgras{šaukštas} \textsuperscript{(34)},
\VUgras{šakutė} \textsuperscript{(35)}, \VUgras{peilis} \textsuperscript{(36)} ir
\VUgras{stiklinė} \textsuperscript{(32)}. Buvo ir \VUgras{buteliai}
\textsuperscript{(30)} vynui bei \VUgras{grafinas} \textsuperscript{(31)} vandeniui.

%%K t06-c3-08-1 p
8. --- \VUgras{Tarnas} \textsuperscript{(39)} pirmiausia atnešė \VUgras{sriubinę}
\textsuperscript{(37)}, kurioje dar garuoja truputis sriubos. Dabar jis patiekia
pirmąjį \VUgras{patiekalą} \textsuperscript{(38)}, mėsos patiekalą. Paskui jis
pateiks tuos, kuriuos jau įvardijome, galiausiai, po \VUgras{deserto,} jis
pasinaudos tuo, kad šeimininkai bus išėję į verandą, kad nurinktų stalo reikmenis ir
vėl sutvarkytų valgomąjį.

%%K t06-c3-08-2 p
\textit{Pastaba.} --- \textit{Bendriniai žodžiai, susiję su maistu, čia nurodomi labai
glaustai. Sąrašas bus papildytas dviejose lentelėse: „Viešbutis“ (Nr.~13) ir
„Turgus“ (Nr.~15).}

%%K t06-tit-7 sub
\VUsaut{4.6mm}
\VUpk{10.2pt}{40}{Virtuvė.}
\VUsaut{3.0mm}

%%K t06-c4-01-1 p
1. --- Virtuvėje, į kurią žvelgiame pro praviras duris, matome vaizdingą netvarką.
Prieš \VUgras{židinį} \textsuperscript{(1)}, kuriame traška \VUgras{ugnis}
\textsuperscript{(3)}, įtaisytas \VUgras{iešmo sukiklis} \textsuperscript{(5)}. Gražus
\VUgras{kepsnys} \textsuperscript{(7)} \VUgras{užmautas ant iešmo}
\textsuperscript{(6)} ir kepa.

%%K t06-c4-02-1 p
2. --- \VUgras{Dumplės} \textsuperscript{(8)} ir \VUgras{anglių semtuvas}
\textsuperscript{(9)}, kuriais ką tik naudotasi, liko ant grindų, kaip ir
\VUgras{malkos} \textsuperscript{(10)}.

%%K t06-c4-03-1 p
3. --- Židinio viršus sutvarkytas: \VUgras{žibintas} \textsuperscript{(11)},
\VUgras{degtukinė} \textsuperscript{(13)}, kurioje yra \VUgras{degtukai}
\textsuperscript{(12)}, \VUgras{žvakidė} \textsuperscript{(14)} ir
\VUgras{žvakidėlė} \textsuperscript{(15)} su \VUgras{žvake} \textsuperscript{(16)}
yra savo vietoje. Bet didelis \VUgras{anglių kibiras} \textsuperscript{(18)}, pilnas
\VUgras{akmens anglių} \textsuperscript{(17)}, kurių \VUgras{virėja}
\textsuperscript{(23)} ką tik prisipylė rūsyje, liko virtuvės viduryje.

%%K t06-c4-04-1 p
4. --- Iš tiesų vargšei moteriai reikia skubėti, kad pusryčiai būtų paruošti laiku.
Dabar ji prie \VUgras{viryklės} \textsuperscript{(19)} ir maišo \VUgras{padažą}
\textsuperscript{(24)}, kurio nedera per daug prideginti.

%%K t06-c4-05-1 p
5. --- \VUgras{Orkaitėje} \textsuperscript{(20)} kepa pyragas, kurį ji paruošė
vaikų užkandžiui. Virš viryklės kabo \VUgras{samtis} \textsuperscript{(21)} ir
\VUgras{putų graibštas} \textsuperscript{(22)}.

%%K t06-tit-8 sub
\VUsaut{4.0mm}
\VUpk{10.2pt}{40}{Indai.}
\VUsaut{3.0mm}

%%K t06-c4-06-1 p
6. --- \VUgras{Virtuvės indai} \textsuperscript{(25)}, sukabinti kambario gale,
labai švarūs. Gražūs variniai \VUgras{prikaistuviai} \textsuperscript{(26)},
\VUgras{pieninė} \textsuperscript{(27)}, \VUgras{sietelis} \textsuperscript{(28)} ir
\VUgras{tarka} \textsuperscript{(29)} buvo rūpestingai nuvalyti.

%%K t06-c4-07-1 p
7. --- \VUgras{Druskinė} \textsuperscript{(30)} padėta ant indų spintos prie
\VUgras{arbatinuko} \textsuperscript{(31)} ir \VUgras{aliejaus butelio}
\textsuperscript{(32)}. \VUgras{Stalčiuje} \textsuperscript{(33)} tikriausiai laikomi
stalo įrankiai ir virtuvės peiliai.

%%K t06-c4-08-1 p
8. --- \VUgras{Šluota} \textsuperscript{(34)} ir \VUgras{plunksnų šluostiklis}
\textsuperscript{(35)} kampe. Virėja ką tik naudojosi \VUgras{medine trinka}
\textsuperscript{(36)}, ji paliko ją prie lango.

%%K t06-c4-09-1 p
9. --- \VUgras{Rankšluostis} \textsuperscript{(37)} kabo ant sienos, prie
\VUgras{piltuvo} \textsuperscript{(38)}. Toje virtuvės dalyje padėtos
\VUgras{grotelės} \textsuperscript{(39)}, \VUgras{kapoklė} \textsuperscript{(40)} ir
\VUgras{kapojimo lentelė} \textsuperscript{(41)}. Paskui, virš kapoklės, ant
\VUgras{lentynos} \textsuperscript{(42)}, yra \VUgras{puodai}
\textsuperscript{(43)}, \VUgras{virdulys} \textsuperscript{(44)} su savo
\VUgras{dangčiu} \textsuperscript{(45)} ir \VUgras{katilas} \textsuperscript{(46)}.
\VUgras{Keptuvė} \textsuperscript{(47)} kabo netoliese.

%%K t06-c4-10-1 p
10. --- Tik varinis \VUgras{čiaupas} \textsuperscript{(48)} meta blizgesį šiame
kampe. \VUgras{Plautuvė} \textsuperscript{(49)}, ant kurios pastatytas storas
\VUgras{ąsotis} \textsuperscript{(50)}, tikriausiai labai patogi su savo mažute
spintele, kurioje laikoma \VUgras{acto statinaitė} \textsuperscript{(51)}, aprūpinta
savo \VUgras{čiaupeliu} \textsuperscript{(52)}, \VUgras{atliekų dėžė}
\textsuperscript{(53)} ir \VUgras{nešvarūs indai} \textsuperscript{(54)}.

%%K t06-tit-9 sub
\VUsaut{4.0mm}
\VUpk{10.2pt}{40}{Kiti virtuvėje padėti daiktai.}
\VUsaut{3.0mm}

%%K t06-c4-11-1 p
11. --- \VUgras{Kubilas} \textsuperscript{(55)}, kuriame plaunamos
\VUgras{daržovės} \textsuperscript{(58)}, ir ketaus \VUgras{katilas}
\textsuperscript{(56)}, pastatytas ant \VUgras{plytelių grindų}
\textsuperscript{(57)}, tikriausiai irgi turi savo vietą toje spintelėje.

%%K t06-c4-12-1 p
12. --- Virėjai tikrai netrūksta viryklių, nes štai dar, ant stalo kampo,
\VUgras{dujinė viryklė} \textsuperscript{(59)}, ant kurios mažame prikaistuvyje šyla
vanduo. \VUgras{Garai} \textsuperscript{(60)}, kurie iš jo veržiasi, rodo mums, kad
jis tikriausiai verda. Turbūt tas vanduo bus panaudotas kavai, nes štai kitame stalo
gale \VUgras{kavinukas} \textsuperscript{(64)} ir \VUgras{kavamalė}
\textsuperscript{(63)}.

%%K t06-c4-13-1 p
13. --- Viryklė sujungta su \VUgras{dujų degikliu} \textsuperscript{(62)} ilgu
\VUgras{guminiu vamzdeliu} \textsuperscript{(61)}.

%%K t06-c4-14-1 p
14. --- \VUgras{Dieninė tarnaitė} \textsuperscript{(66)}, kuri ateina padėti
virėjai, ruošia daržoves vakarienei. Kai jos bus nuvalytos ir nuplautos, ji sudės jas
į \VUgras{dubenį} \textsuperscript{(65)}, laukdama valandos joms virti.

%%K t06-tit-10 sub
\VUsaut{4.0mm}
{\centering\fontsize{10.2pt}{10.2pt}\selectfont\textls[40]{\scshape Vonia.} \textit{(Maudymosi kambarys.)}\par}
\VUsaut{3.0mm}

%%K t06-c5-01-1 p
1. --- Mums belieka padaryti tik vieną nekuklumą, kad pažintume pagrindinius buto
kambarius: pamatyti vonios įrengimą. Tai nepareikalaus daug laiko, nes ji nedidelė,
ir greitai sužinosime, kad \VUgras{vonia} \textsuperscript{(1)} turi gerą
\VUgras{vandens šildytuvą} \textsuperscript{(2)}, kad \VUgras{muilinė}
\textsuperscript{(3)} pasiekiama besimaudančiojo ranka ir kad \VUgras{rankšluosčių
kabykla} \textsuperscript{(5)} tikrai leis lengvai išdžiovinti \VUgras{rankšluosčius}
\textsuperscript{(4)}.

%%K t06-c5-02-1 p
2. --- To kambario sienos padengtos \VUgras{fajanso plytelėmis}
\textsuperscript{(6)}, kurios leido įrengti \VUgras{dušo įrenginį}
\textsuperscript{(7)} visai nebijant, kad vanduo, kuris tykšta jį naudojant, sugadins
sienas. Cinkinis \VUgras{dubenėlis} \textsuperscript{(8)}, kuris skirtas kojų
vonelėms, papildo įrangą.

\VUsaut{6.0mm}
\VUfilet{22mm}
